\section{Scope \& Methodology}
\subsection{In-Scope Networks \& Sites}
\begin{center}\begin{tabularx}{\linewidth}{@{}l X@{}}
\toprule\rowcolor{TableHead}\textcolor{white}{\bfseries Site} & \textcolor{white}{\bfseries Networks} \\
\midrule
\fillme{Building/floor} & \fillme{SSIDs in scope} \\
\bottomrule\end{tabularx}\end{center}

\subsection{Rules of Engagement}
\begin{itemize}
  \item Testing window: \fillme{start date} -- \fillme{end date}, on-site.
  \item Deauthentication / jamming attacks: \fillme{in / out of scope}.
  \item Rogue AP impersonation: \fillme{in / out of scope}.
\end{itemize}

\subsection{Methodology}
\begin{enumerate}
  \item \textbf{RF survey.} Passive scanning to enumerate all in-range access points, SSIDs, and
        encryption types (\texttt{airodump-ng}/\texttt{Kismet}).
  \item \textbf{Encryption assessment.} Handshake/PMKID capture and offline cracking attempts
        against WPA2-PSK networks; EAP-method identification for WPA2/3-Enterprise networks.
  \item \textbf{Rogue AP / Evil Twin.} Deploying a cloned SSID to test client auto-connect
        behaviour and credential capture (\texttt{hostapd-wpe} or similar).
  \item \textbf{Client isolation \& segmentation testing.} Verifying guest networks cannot reach
        internal/corporate segments.
  \item \textbf{Captive portal testing.} Bypass attempts against any captive-portal authentication.
\end{enumerate}

\subsection{Severity Definitions}
\begin{center}\begin{tabularx}{\linewidth}{@{}l X@{}}
\toprule\rowcolor{TableHead}\textcolor{white}{\bfseries Level} & \textcolor{white}{\bfseries Definition} \\
\midrule
\critical & Corporate credential capture or direct access to internal network from wireless. \\
\rowcolor{TableAlt}
\high     & Crackable encryption, weak EAP method, or captive-portal bypass. \\
\medium   & Missing client isolation, weak guest-network segmentation. \\
\rowcolor{TableAlt}
\low\ / Info & Hygiene gaps (e.g. SSID broadcasting internal naming conventions). \\
\bottomrule\end{tabularx}\end{center}
